\documentclass[11pt]{article}
\usepackage[a4paper,margin=1in]{geometry}
\usepackage{amsmath,amssymb,amsthm}
\usepackage[T1]{fontenc}
\usepackage{lmodern}
\usepackage[hidelinks]{hyperref}
\setlength{\emergencystretch}{3em}

\theoremstyle{plain}
\newtheorem{theorem}{Theorem}
\newtheorem{lemma}[theorem]{Lemma}
\theoremstyle{definition}
\newtheorem{definition}[theorem]{Definition}
\theoremstyle{remark}
\newtheorem{remark}[theorem]{Remark}

\newcommand{\bx}{\operatorname{box}}
\newcommand{\R}{\mathbb R}

\title{A Counterexample to Conjecture 00000003965:\\
Boxicity of $K_{2,2,2,2}$ and of Its Complement}
\author{%
  Feiyu\thanks{Independent researcher.}
  \and
  Claude (Opus 5.5)\thanks{Anthropic.}%
}
\date{2026-10-03}

\begin{document}
\maketitle

\begin{abstract}
Conjecture 00000003965 asserts that every graph $G$ satisfies
$|\bx(G)-\bx(\overline G)|\le2$. For the complete multipartite graph
$G=K_{2,2,2,2}$ we have $\bx(G)=4$, while its complement is a perfect matching with
$\bx(\overline G)=1$. So the difference is $3$. The refutation is checked in Lean~4
against Mathlib.
\end{abstract}

\section{The conjecture as stated}

The original text of conjecture 00000003965 reads:

\begin{quote}
Definition: The boxicity of a graph is the minimal dimension in which the graph is
represented as an intersection of interval boxes. Conjecture: Every graph $G$
satisfies $|\bx(G) - \bx(\text{complement of } G)| \le 2$ with respect to its
complement, and the difference $2$ is attained by an explicit construction.
\end{quote}

The Chinese version in \texttt{conjectures/00000003965.md} states the same claim.

\section{Reading of the statement}

\begin{definition}
A \emph{$d$-dimensional box representation} of a graph $G$ assigns to every vertex
$v$ an axis-parallel box $\prod_{i=1}^d[l_i(v),r_i(v)]\subseteq\R^d$, with
$l_i(v)\le r_i(v)$, such that two distinct vertices are adjacent if and only if their
boxes intersect. Two boxes intersect iff their intervals overlap in every coordinate,
where $[a,b]$ and $[c,e]$ overlap iff $a\le e$ and $c\le b$. The \emph{boxicity}
$\bx(G)$ is the least $d$ for which such a representation exists.
\end{definition}

In dimension $0$ every box is the single point $\R^0$, so exactly the complete graphs
have boxicity $0$; this is the usual convention. The first claim of the conjecture is
$|\bx(G)-\bx(\overline G)|\le2$ for every graph $G$, where $\overline G$ is the
complement. We refute it with a finite graph, so restricting the claim to finite
graphs makes no difference. We do not need the second claim.

\section{Result}

Let $G=K_{2,2,2,2}$ be the complete multipartite graph with four parts of size two.
Its vertices are the pairs $(i,s)$ with $i\in\{1,2,3,4\}$ and $s\in\{0,1\}$, and
$(i,s)$ is adjacent to $(k,t)$ iff $i\ne k$. The complement $\overline G$ is the
perfect matching $4K_2$ with edges $\{(i,0),(i,1)\}$.

\begin{theorem}\label{thm:main}
$\bx(G)=4$ and $\bx(\overline G)\le1$. Hence $\bx(G)-\bx(\overline G)\ge3$, and
Conjecture 00000003965 is false.
\end{theorem}

This is the case $k=4$ of Roberts' theorem that the complete multipartite graph with
$k$ parts of size two has boxicity $k$~\cite{roberts}. We give the short proof.

\section{Proof}

\begin{lemma}\label{lem:interval}
Let $A,B,C,E$ be closed intervals of $\R$ such that $A\cap B=\emptyset$ and each of
$C$ and $E$ meets both $A$ and $B$. Then $C\cap E\ne\emptyset$.
\end{lemma}

\begin{proof}
Say $A=[a_1,a_2]$ lies to the left of $B=[b_1,b_2]$, so $a_2<b_1$. Since $C=[c_1,c_2]$
meets $A$ and $B$, we have $c_1\le a_2$ and $c_2\ge b_1>a_2$, so $a_2\in C$. Likewise
$a_2\in E$.
\end{proof}

\begin{proof}[Proof of Theorem~\ref{thm:main}]
\emph{Upper bound for $\overline G$.} Give $(i,0)$ and $(i,1)$ the interval
$[3i,3i+1]$. Intervals of the same $i$ coincide, and intervals of different $i$ are
disjoint, so this is a $1$-dimensional representation of $\overline G$.

\emph{Upper bound for $G$.} In coordinate $j\in\{1,2,3,4\}$ give $(j,0)$ the interval
$[0,1]$, give $(j,1)$ the interval $[2,3]$, and give every vertex $(i,s)$ with
$i\ne j$ the interval $[0,3]$. In coordinate $j$ the only non-overlapping pair is
$\{(j,0),(j,1)\}$. Hence two distinct vertices have intersecting boxes iff they do not
form one of the pairs $\{(j,0),(j,1)\}$, that is, iff they are adjacent in $G$. So
$\bx(G)\le4$.

\emph{Lower bound for $G$.} Suppose $G$ has a $d$-dimensional representation with
$d\le3$. The vertices $(i,0)$ and $(i,1)$ are distinct and non-adjacent, so their
boxes are disjoint and there is a coordinate $\mathrm{sep}(i)$ in which their
intervals are disjoint. Since $\mathrm{sep}$ maps four values into $d\le3$
coordinates, $\mathrm{sep}(i)=\mathrm{sep}(k)=j$ for some $i\ne k$. In coordinate $j$
the intervals $A,B$ of $(i,0),(i,1)$ are disjoint. The intervals $C,E$ of
$(k,0),(k,1)$ each meet $A$ and $B$, because $(k,t)$ is adjacent to $(i,s)$. By
Lemma~\ref{lem:interval}, $C$ and $E$ meet, contradicting the choice of
$j=\mathrm{sep}(k)$. Hence $\bx(G)\ge4$.
\end{proof}

\section{Formalization}

\sloppy
The Lean~4 project in \texttt{lean/} (file \texttt{lean/Submission/Basic.lean})
formalizes the definitions above and proves Theorem~\ref{thm:main} against Mathlib.
\begin{itemize}
  \item \texttt{BoxRep G d} is a $d$-dimensional box representation: real endpoint
    maps $l,r$ on vertices and coordinates $i\in\texttt{Fin d}$ with
    $l_i(v)\le r_i(v)$, such that distinct $u,v$ are adjacent iff
    $l_i(u)\le r_i(v)$ and $l_i(v)\le r_i(u)$ for all $i$. \texttt{boxicity G} is the
    infimum (\texttt{sInf}) of the dimensions admitting one.
  \item \texttt{BoxicityGapAtMostTwo} states $|\bx(G)-\bx(\overline G)|\le2$ for
    every finite graph, with $\overline G$ Mathlib's complement of \texttt{G}.
  \item The graph $G$ is Mathlib's \texttt{completeEquipartiteGraph 4 2} on
    \texttt{Fin 4} $\times$ \texttt{Fin 2}.
  \item \texttt{repCompl} and \texttt{repK} are the two representations above; the
    adjacency conditions are checked by \texttt{decide} on integer endpoints.
  \item \texttt{interval\_lemma} is Lemma~\ref{lem:interval}, and
    \texttt{no\_rep\_of\_lt\_four} is the lower bound, using the pigeonhole principle
    \texttt{Fintype.exists\_ne\_map\_eq\_of\_card\_lt}.
  \item \texttt{boxicity\_K} gives $\bx(G)=4$, \texttt{boxicity\_compl\_le} gives
    $\bx(\overline G)\le1$, and \texttt{conjecture\_00000003965\_false} is
    $\lnot\,$\texttt{BoxicityGapAtMostTwo}.
\end{itemize}
The toolchain is \texttt{leanprover/lean4:v4.33.1}, Mathlib is pinned to
\texttt{v4.33.1}, and \texttt{lake build} succeeds. The project contains no
\texttt{sorry}, no \texttt{admit} and no \texttt{native\_decide}, and introduces no
axioms: \texttt{\#print axioms} reports exactly \texttt{propext},
\texttt{Classical.choice} and \texttt{Quot.sound} for
\texttt{conjecture\_00000003965\_false} and \texttt{boxicity\_K}.

\fussy
\section{Remarks}

\begin{remark}
The same construction gives unbounded differences: for the complete multipartite
graph with $k$ parts of size two, $\bx=k$ while the complement (a perfect matching)
has boxicity $1$~\cite{roberts}. So $|\bx(G)-\bx(\overline G)|$ is not bounded by any
constant.
\end{remark}

\begin{thebibliography}{9}
\bibitem{roberts} F.~S.~Roberts, \emph{On the boxicity and cubicity of a graph}, in:
  Recent Progress in Combinatorics (W.~T.~Tutte, ed.), Academic Press, 1969,
  301--310.
\end{thebibliography}

\end{document}
